%%%%%%%%%%%%%%%%%%%%%%%%%%%%%%%%%%%%%%%%%%%%%%%%%%%%%%%%%%%%%%%%%%%%%%%%%%%%%%%%
%
%   大连理工大学硕士论文模板 ―― 主文件 main.tex
%   版本：1.0
%   最后更新：2024-01-01
%   编译环境：XeLaTeX + biblatex + biber
%   欢迎使用，如有任何意见和建议请发送邮件至 yuri_1985@163.com，感谢您的支持！
%
%%%%%%%%%%%%%%%%%%%%%%%%%%%%%%%%%%%%%%%%%%%%%%%%%%%%%%%%%%%%%%%%%%%%%%%%%%%%%%%%

\documentclass[UTF8]{dlutthesis}

% 定义所有的图片文件在 figures 子目录下
\graphicspath{{figures/}}

% 参考文献数据库
\addbibresource{sections/reference.bib}

% 论文信息设置
\ctitle{大连理工大学学位论文 \\ XeLaTeX 模板设计与实现}
\cdegree{硕士学位论文}
\cauthor{作者姓名}
\cauthorno{2020000000}
\csupervisor{导师姓名\quad 教授}
\csubject{专业名称}
\cschool{学院（学部）名称}
\creviewer{评阅教师}
\cdate{2024年6月}
\cabstract{
    这里是中文摘要内容。本论文主要介绍了大连理工大学学位论文XeLaTeX模板的设计与实现过程。
    通过分析现有模板的不足，结合现代LaTeX技术，重新设计了一个基于.cls文档类的模板，
    提高了模板的易用性、可维护性和编译效率。

    本文首先回顾了LaTeX在学术写作中的应用现状，然后详细阐述了模板的设计思路和实现细节，
    包括文档类结构、格式设置、宏包集成等方面。最后，通过对比测试验证了新模板的优越性。
}
\ckeywords{XeLaTeX；学位论文模板；文档类；自动化排版}

\etitle{Design and Implementation of XeLaTeX Template for Dalian University of Technology}
\eabstract{
    This is the English abstract. This thesis mainly introduces the design and implementation
    of the XeLaTeX template for Dalian University of Technology degree thesis. Through analyzing
    the shortcomings of existing templates and combining modern LaTeX technology, a new template
    based on .cls document class is redesigned to improve the usability, maintainability and
    compilation efficiency of the template.

    This paper first reviews the current application status of LaTeX in academic writing,
    then elaborates the design ideas and implementation details of the template, including
    document class structure, format settings, package integration, etc. Finally, the superiority
    of the new template is verified through comparative testing.
}
\ekeywords{XeLaTeX; thesis template; document class; automated typesetting}

% 文档正文
\begin{document}

% 前言
\frontmatter
\pagenumbering{Roman}

% 封面
\makecover

% 独创性声明
% \originality

% 摘要
% \makeabstract

% 设置目录字体和行间距
\defaultmenufont

% 目录
\tableofcontents
\cleardoublepage

% 插图目录
%\listoffigures

% 表格目录
%\listoftables

\cleardoublepage
\defaultfont
\mainmatter

% 正文章节
\defaultfont
\renewcommand{\thefootnote}{\arabic{footnote}}
\chapter*{\hfill 引　　言 \hfill}
\addcontentsline{toc}{chapter}{引　　言}
\label{chap00:introduction}
\chapter{git仓库分析}
\label{chap01}
\section{基于gitgraph的git提交历史可视化}
本章详细介绍如何使用 GitGraph.js 库实现 Git 提交历史的可视化展示。通过分析 Git 仓库的提交数据，生成直观的图形化界面，帮助开发者更好地理解项目的发展历史和分支结构。

\subsection{技术架构设计}

\subsubsection{GitGraph.js 核心特性}

GitGraph.js 是一个基于 HTML5 Canvas 的 JavaScript 库，专门用于绘制 Git 图形。该库具备多分支支持能力，能够准确展示复杂的 Git 分支结构，包括分支的创建、合并、分叉等各种操作历史。库提供了丰富的交互式界面功能，支持用户通过鼠标滚轮进行缩放控制，可以拖拽浏览图形的不同部分，让开发者能够从各个角度详细查看提交历史。在样式配置方面，GitGraph.js 提供了高度自定义的选项，允许开发者根据需求调整颜色、布局、字体等各种视觉元素。此外，该库采用响应式设计原则，能够自动适配不同屏幕尺寸，确保在各种设备上都能提供良好的用户体验。

\subsubsection{系统架构}

整个可视化系统采用清晰的三层架构设计。数据层作为基础层，通过 GitPython 库与 Git 仓库进行深度交互，解析并提取详细的提交信息，包括提交哈希、作者信息、时间戳、父提交关系以及分支标签等核心数据。处理层承担着承上启下的关键角色，Python 脚本不仅对原始 Git 数据进行结构化处理和逻辑分析，还需要将这些数据注入到精心设计的 HTML 模板中，同时集成 GitGraph.js 的配置参数和样式设置。展示层作为最上层，最终通过现代浏览器强大的渲染能力，将抽象的 Git 历史数据转化为直观的可视化图形界面，用户可以清晰地看到分支的创建、合并、分叉等复杂操作的历史轨迹。

\subsection{核心实现}

\subsubsection{HTML 模板生成}

主要实现在 \texttt{html\_generator.py} 文件中，核心函数 \texttt{generate\_git\_tree\_html} 负责生成完整的 HTML 页面：

\begin{verbatim}
def generate_git_tree_html(commits, git_url, output_path="git_tree.html"):
    html_content = f"""
    <!DOCTYPE html>
    <html>
    <head>
        <meta charset="UTF-8">
        <title>Git Tree Visualization</title>
        <script src="https://unpkg.com/@gitgraph/js/lib/gitgraph.umd.js">
        </script>
    </head>
    <body>
        <h1>Git Repository History - {git_url}</h1>
        <div id="graph-container"></div>
    </body>
    </html>
    """
\end{verbatim}

\subsubsection{样式配置}

为了提升用户体验，系统采用 Metro 风格的视觉设计，在色彩方案方面使用 Material Design 色彩体系，提供了 10 种不同的分支颜色，确保各个分支能够清晰区分。在布局优化方面，系统采用垂直反向布局方式，将最新的提交显示在顶部，这样用户可以直观地看到项目的最新进展。字体设置充分考虑了中英文混合显示的需求，优先使用 Consolas 和微软雅黑字体，既保证了代码的可读性，又确保了中文字符的美观显示。此外，系统还添加了柔和的阴影效果，通过微妙的层次感提升整体视觉体验，让图形界面更加立体和现代化。

\subsubsection{分支管理算法}

系统实现了智能的分支管理算法，能够准确处理以下场景：

\paragraph{主干分支处理}

\begin{verbatim}
const main = gitgraph.branch({
    name: "main",
    style: { label: { display: true } }
});
branches["main"] = main;
branchToLastCommit.set(main, null);
\end{verbatim}

\paragraph{分支创建与切换}

当检测到新的分支引用时，系统会自动创建对应的分支：

\begin{verbatim}
c.refs.forEach(ref => {
    if (ref.startsWith('refs/heads/') ||
        (ref.startsWith('origin/') && !ref.includes('HEAD'))) {
        const bName = ref.replace('refs/heads/', '')
                        .replace('origin/', '');
        if (!branches[bName]) {
            branches[bName] = targetBranch.branch({
                name: bName,
                style: { label: { display: true } }
            });
        }
    }
});
\end{verbatim}

\paragraph{合并操作处理}

对于合并提交（包含多个父提交），系统会智能识别并执行合并操作：

\begin{verbatim}
if (c.parents.length > 1) {
    for (let i = 1; i < c.parents.length; i++) {
        const otherBranch = commitToBranch[c.parents[i]];
        if (otherBranch && otherBranch !== targetBranch) {
            targetBranch.merge({
                branch: otherBranch,
                commitOptions: commitOptions
            });
            break;
        }
    }
}
\end{verbatim}

\subsection{功能特性}

\subsubsection{提交信息展示}

每个提交节点都包含丰富的信息展示，在提交哈希方面显示完整的 SHA-1 哈希值，方便开发者进行精确的定位和查找。提交信息部分则显示详细的提交说明文字，让开发者能够快速了解每次提交的具体内容和目的。作者信息包含了作者姓名和提交时间，为团队协作提供了必要的时间线和责任人信息。在视觉标识方面，系统采用圆点表示提交节点，并通过不同颜色区分各个分支，使得复杂的 Git 历史结构一目了然。

\subsubsection{交互功能}

系统提供了丰富的交互功能，在缩放控制方面支持鼠标滚轮缩放图形，用户可以根据需要放大查看细节或缩小查看整体结构。拖拽浏览功能允许用户拖拽查看图形的不同部分，即使对于包含大量提交历史的项目也能轻松导航。分支标签功能清晰显示各分支名称，帮助用户快速识别和定位不同的开发线路。响应式适配机制确保图形能够自动适应容器大小，无论是在大型显示器还是小型设备上都能够提供最佳的视觉效果。

\subsection{性能优化}

\subsubsection{数据处理优化}

在数据处理方面，系统采用批量处理策略，一次性处理所有提交数据，避免了重复计算带来的性能损耗。内存管理方面使用 Map 数据结构来优化查找性能，确保在处理大量提交数据时仍能保持高效的查询速度。引用缓存机制则通过缓存分支和提交的映射关系，减少了重复的查找操作，显著提升了整体处理效率。

\subsection{输出结果}

生成的 HTML 文件包含完整的 Git 图形可视化，用户可以在浏览器中直接打开查看，无需额外的配置或依赖。该文件可以作为静态文件保存，供后续重复使用或分享给团队成员。同时，生成的 HTML 代码也可以轻松嵌入到其他网页中，作为项目文档或报告的一部分。此外，系统还支持打印和截图功能，方便用户将可视化结果保存为图片或 PDF 格式，用于演示文稿或文档编制。图\ref{fig:git_history}展示了生成结果的可视化效果。

\begin{figure}[htbp]
\centering
\includegraphics[width=0.9\textwidth]{figures/git_history.png}
\caption{Git 提交历史可视化效果图}
\label{fig:git_history}
\end{figure}


\chapter{代码静态/动态分析}
\label{chap02}

\chapter{注意事项}
\label{chap03}


% 结论
\chapter*{\hfill 结　　论 \hfill}
\addcontentsline{toc}{chapter}{结　　论}

本文设计并实现了一套多维度的软件仓库分析工具，通过整合 Git 版本控制数据、代码静态分析以及 GitHub 平台元数据，为软件项目的质量评估与演进分析提供了全面的解决方案。

主要工作总结如下：

\begin{enumerate}
    \item \textbf{可视化 Git 历史分析}：通过 \texttt{gitgraph.js} 和多维度统计图表，直观展示了项目的开发演进路线、团队贡献分布及开发节奏，帮助管理者快速把握项目现状。
    \item \textbf{深度代码质量评估}：集成了圈复杂度计算、依赖树分析及漏洞扫描功能。能够自动识别高复杂度代码块、过时或危险的第三方依赖，以及潜在的代码安全漏洞，为代码重构和安全加固提供精准指引。
    \item \textbf{GitHub 协作流程洞察}：实现了对 Issue, Pull Request 及 GitHub Actions 的自动化分析。不仅量化了社区的响应效率（如 PR 处理周期），还通过关键词挖掘识别了潜在的安全风险，并评估了 CI/CD 流程的效率与安全性。
\end{enumerate}

本工具采用模块化设计，易于扩展。未来的工作将集中在支持更多编程语言的分析、引入 AI 模型进行更深层次的代码语义分析，以及开发实时的仪表盘监控系统，进一步提升工具的实用性与智能化水平。


\cleardoublepage

% 参考文献（使用 biblatex）
\wuhao
\phantomsection
\addcontentsline{toc}{chapter}{参考文献}
\printbibliography[title=参考文献]
\cleardoublepage

%附录
\defaultfont
\begin{appendix}
    \chapter*{\hfill 附录~A　原始分析数据摘录 \hfill}
\addcontentsline{toc}{chapter}{附录~A　原始分析数据摘录}

为了佐证正文中的分析结论，本附录摘录了部分自动化工具生成的原始分析报告。这些数据直接来源于工具链的输出，未进行人工修饰。

\section*{A.1 Bandit 安全扫描报告（节选）}

以下是 \texttt{bandit} 工具对源代码进行 AST 静态扫描后输出的部分高风险与中风险问题记录。报告详细列出了漏洞类型、严重程度及对应的代码位置。

\begin{verbatim}
Run started:2026-02-09 03:44:55.548818+00:00

Test results:
>> Issue: [B606:start_process_with_no_shell] Starting a process without a shell.
   Severity: Low   Confidence: Medium
   CWE: CWE-78 (https://cwe.mitre.org/data/definitions/78.html)
   Location: repo/COMTool/Main.py:16:4
   15          log.i("Restarting..., comand: {} {} {}".format(python, python, *sys.argv))
   16          os.execl(python, python, * sys.argv)

--------------------------------------------------
>> Issue: [B113:request_without_timeout] Call to requests without timeout
   Severity: Medium   Confidence: Low
   CWE: CWE-400 (https://cwe.mitre.org/data/definitions/400.html)
   Location: repo/COMTool/autoUpdate.py:25:19
   24              try:
   25                  page = requests.get(self.releaseApiUrl)
\end{verbatim}

\section*{A.2 Radon 圈复杂度分析（节选）}

以下是 \texttt{radon} 工具生成的圈复杂度 (Cyclomatic Complexity) 原始数据。每行显示了文件路径、类型（M=Method, C=Class, F=Function）、代码行号、名称以及复杂度评分（A为最优，F为最差）。

\begin{verbatim}
repo/COMTool/Combobox.py
    M 23:4 ComboBox._showPopup - A (3)
    C 5:0 ComboBox - A (2)
    M 9:4 ComboBox.__init__ - A (1)
    M 15:4 ComboBox.mouseReleaseEvent - A (1)

repo/COMTool/Main.py
    F 10:0 restart_program - A (1)

repo/COMTool/autoUpdate.py
    M 21:4 AutoUpdate.checkUpdate_github - B (7)
    C 10:0 AutoUpdate - A (4)
    M 50:4 AutoUpdate.checkUpdate_neucrack - A (4)
    M 14:4 AutoUpdate.detectNewVersion - A (2)
    M 72:4 AutoUpdate.decodeTag - A (2)
\end{verbatim}

\section*{A.3 Pylint 代码风格检查（节选）}

以下是 \texttt{pylint} 对代码规范性检查的输出样本。主要包含了命名规范（Snake Case）、文档字符串缺失（Missing docstring）以及导入错误等问题。

\begin{verbatim}
************* Module COMTool
repo/COMTool/__init__.py:1:0: C0103: Module name "COMTool" doesn't conform to 
snake_case naming style (invalid-name)

************* Module COMTool.Combobox
repo/COMTool/Combobox.py:1:0: C0103: Module name "Combobox" doesn't conform to 
snake_case naming style (invalid-name)
repo/COMTool/Combobox.py:5:0: C0115: Missing class docstring (missing-class-docstring)
repo/COMTool/Combobox.py:11:8: C0103: Variable name "listView" doesn't conform to 
snake_case naming style (invalid-name)
repo/COMTool/Combobox.py:15:4: C0103: Method name "mouseReleaseEvent" doesn't conform 
to snake_case naming style (invalid-name)
repo/COMTool/Combobox.py:18:4: C0103: Method name "showPopup" doesn't conform to 
snake_case naming style (invalid-name)
\end{verbatim}

\section*{A.4 PR分析报告}

本节展示了针对 Neutree/COMTool 仓库的 PR 深度分析报告。

\textbf{仓库名称}：Neutree/COMTool \quad
\textbf{分析时间}：2026-02-09 11:46:32 \quad
\textbf{僵尸PR阈值}：7天 \quad
\textbf{异常值过滤阈值}：30天

\subsection*{核心指标汇总}
\begin{longtable}{l|l}
\hline
\textbf{指标} & \textbf{数值} \\
\hline
\endfirsthead
\hline
\textbf{指标} & \textbf{数值} \\
\hline
\endhead
\hline
\endfoot
总PR数 & 23 \\
合入PR数 & 16 \\
拒绝PR数 & 5 \\
未处理PR数 & 2 \\
PR合入率 & 69.57\% \\
PR拒绝率 & 21.74\% \\
僵尸PR数 & 2 \\
僵尸PR占比 & 8.7\% \\
平均审核时长（小时） & 7.86 \\
中位数审核时长（小时） & 1.58 \\
过滤异常值后审核时长 & 7.86 \\
平均评论数 & 1.65 \\
平均提交次数 & 1.57 \\
平均新增代码 & 138.78 行 \\
平均删除代码 & 73.78 行 \\
平均修改文件数 & 2.65 个 \\
平均生命周期（天） & 0.9 \\
\hline
\end{longtable}

\subsection*{关键结论}

\begin{enumerate}
    \item \textbf{PR处理效率}：仓库PR合入率为69.57\%，拒绝率为21.74\%，反映代码评审的严格程度。正常的拒绝率可能表明评审标准较为严格。
    \item \textbf{僵尸PR情况}：僵尸PR占比8.7\%，正常的比例表明需要关注PR处理效率。
    \item \textbf{审核效率}：平均审核时长7.86小时，中位数为1.58小时。过滤异常值后为7.86小时，正常的时长可能影响开发效率。
    \item \textbf{代码质量}：平均每次PR新增代码138.78行，删除代码73.78行，修改文件2.65个。较小规模的修改可能增加评审难度。
    \item \textbf{PR生命周期}：关闭PR的平均生命周期为0.9天，正常的生命周期可能表明评审流程需要优化。
    \item \textbf{协作模式}：PR的平均评论数为1.65条，较少的评论表明团队协作程度较高。
    \item \textbf{提交习惯}：平均每次PR提交次数为1.57次，正常的提交次数可能表明开发过程较为迭代。
\end{enumerate}

\subsection*{详细分析}

\subsubsection*{PR作者分布}
\begin{longtable}{l|c|r}
\hline
\textbf{作者} & \textbf{PR数量} & \textbf{占比} \\
\hline
YI-- & 2 & 8.7\% \\
junchao98 & 2 & 8.7\% \\
e71828 & 2 & 8.7\% \\
lhdjply & 2 & 8.7\% \\
taotieren & 2 & 8.7\% \\
dependabot[bot] & 2 & 8.7\% \\
taorye & 2 & 8.7\% \\
wirano & 1 & 4.3\% \\
sunzigang & 1 & 4.3\% \\
kaidegit & 1 & 4.3\% \\
\hline
\end{longtable}

\subsubsection*{PR标签分布}
\begin{longtable}{l|c|r}
\hline
\textbf{标签} & \textbf{PR数量} & \textbf{占比} \\
\hline
无标签 & 21 & 91.3\% \\
dependencies & 2 & 8.7\% \\
\hline
\end{longtable}

\subsubsection*{PR分支分析}

\textbf{目标分支分布}
\begin{longtable}{l|c}
\hline
\textbf{分支} & \textbf{PR数量} \\
\hline
master & 23 \\
\hline
\end{longtable}

\textbf{来源分支分布}
\begin{longtable}{l|c}
\hline
\textbf{分支} & \textbf{PR数量} \\
\hline
master & 13 \\
aur & 2 \\
feat/record & 1 \\
fix/protocols & 1 \\
makensis & 1 \\
develop/savetimesamp & 1 \\
GImDX-patch-1 & 1 \\
fix-badge & 1 \\
dependabot/pip/urllib3-1.26.5 & 1 \\
dependabot/pip/urllib3-1.25.8 & 1 \\
\hline
\end{longtable}

\subsubsection*{PR生命周期分布}
\begin{longtable}{l|c}
\hline
\textbf{生命周期} & \textbf{PR数量} \\
\hline
$<$1天 & 19 \\
1-3天 & 1 \\
3-7天 & 0 \\
7-14天 & 0 \\
14-30天 & 1 \\
30-60天 & 0 \\
$>$60天 & 0 \\
\hline
\end{longtable}

\vspace{1em}
\textit{报告生成时间：2026-02-09 11:46:32}

\end{appendix}
\cleardoublepage

\defaultfont
%发表的文章列表
%\include{sections/8.publications}
\cleardoublepage

%致谢
% \include{sections/7.acknowledgements}
% \cleardoublepage

%授权书
%\chapter*{}
%\addcontentsline{toc}{chapter}{大连理工大学学位论文版权使用授权书}
%\renewcommand{\baselinestretch}{1.61}
%\vspace{-0.48cm}
%\authorization
%\cleardoublepage

\end{document}

%%%%%%%%%%%%%%%%%% End of the file  %%%%%%%%%%%%%%%%%%%%%%%%
